\documentclass[11pt]{article}
\usepackage[margin=1in]{geometry}
\usepackage[T1]{fontenc}
\usepackage[utf8]{inputenc}
\usepackage{lmodern}
\usepackage{microtype}
\usepackage{booktabs}
\usepackage{tabularx}
\usepackage{array}
\usepackage{longtable}
\usepackage{amsmath}
\usepackage{enumitem}
\usepackage[hidelinks]{hyperref}
\newcolumntype{Y}{>{\raggedright\arraybackslash}X}
\newcommand{\chvd}{\textsc{CHVD}}
\newcommand{\mone}{M1}
\newcommand{\mthree}{M3}

\title{\textbf{Controlled Distillation of Speculative LLM Venture Concepts:}\\
Evaluator-Assessed Evidence for Bounded MVP Hypotheses}
\author{Anonymous author(s) for review}
\date{}

\begin{document}
\maketitle

\begin{abstract}
Large language models can generate unusual venture concepts, but speculative outputs
often combine a potentially distinctive mechanism with unsupported claims, broad
dependencies, and an MVP scope that is difficult to test. We introduce Controlled
Hallucination Venture Discovery (\chvd{}), a controlled concept-distillation procedure
that separates speculative venture-concept generation (\mone{}) from a bounded
transformation (\mthree{}). The objective is not to establish market success or
real-world startup readiness. Instead, the study tests whether a transformation can
retain a concept's distinctive creative mechanism while improving evaluator-assessed
prototype readiness at the level of written concept cards. The primary study uses 60
matched \mone{}--\mthree{} pairs from Gemini 3.1 Flash Lite P1 across six domains.
Two independent LLM evaluators performed blinded individual-card scoring, and two
role-aware directional audits assessed creative-core preservation and practical MVP
improvement. In individual scoring, technical plausibility was higher for \mthree{}
than for \mone{} for both evaluators after Holm correction (ChatGPT:
$\Delta=+0.150$, $p_{\mathrm{Holm}}=.025$; Gemini: $\Delta=+0.433$,
$p_{\mathrm{Holm}}<.001$). Directional audits rated creative-core preservation near
4/5, practical MVP improvement above 4.27/5, and overclaim-removal quality above
4/5 for both evaluators. The results support a bounded claim: \chvd{} can function
as an auditable procedure for converting selected speculative mechanisms into more
testable MVP hypotheses. They do not demonstrate human preference, factual
correctness, market viability, startup success, or universal superiority of
transformed concepts.
\end{abstract}

\noindent\textbf{Keywords:} large language models; venture ideation; concept
distillation; hallucination; prototype readiness; LLM-as-a-judge; minimum viable product

\section{Introduction}
Large language models (LLMs) can generate plausible and unusual venture concepts. The
same capacity for non-obvious association can also introduce unsupported technical
claims, unclear dependencies, and scopes too broad to test. In early ideation, the
question is not whether unsupported content is intrinsically valuable. It is whether a
controlled process can retain a distinctive creative mechanism while converting the
mechanism into a constrained experimental hypothesis.

We introduce Controlled Hallucination Venture Discovery (\chvd{}), a two-stage
within-model concept-distillation framework. \mone{} is a speculative venture concept
that may contain a distinctive but unverified mechanism. \mthree{} is a controlled,
bounded reformulation intended to preserve that mechanism while replacing open-ended
dependencies or unsupported assumptions with a testable MVP hypothesis.

The contribution is deliberately narrow. We report evaluator-assessed properties of
written cards rather than market outcomes. We also separate individual-card scoring,
role-aware directional traceability, and neutral paired comparison because those
protocols answer different questions and cannot be treated as interchangeable
evidence.

\section{Study design and evidence policy}
Table~\ref{tab:inventory} distinguishes the primary Gemini 3.1 Flash Lite P1 study
from earlier Gemini 2.5 Flash Lite P0/P1 developmental studies. The studies differ in
corpus construction, domains, prompt regimes, and evaluation sessions; raw scores are
not pooled and no generator-model ranking is claimed.

\begin{table}[htbp]\centering\small
\caption{Study inventory and evidential role.}
\label{tab:inventory}
\begin{tabularx}{\textwidth}{@{}lrrYY@{}}
\toprule
Study & Pairs & Role & Individual/directional evidence & Neutral paired status\\
\midrule
Gemini 2.5 FL P0 & 38 & Supporting & Developmental only; not pooled & Not used for main inference\\
Gemini 2.5 FL P1 & 35 & Supporting & Developmental only; not pooled & Non-confirmatory after Holm\\
Gemini 3.1 FL P1 & 60 & Primary & Main individual and directional evidence & Quarantined; not used for efficacy/equivalence\\
\bottomrule
\end{tabularx}
\end{table}

The primary individual protocol compares matched \mone{}--\mthree{} cards for ten
metrics using paired Wilcoxon signed-rank tests, with Holm correction within each
evaluator. The cross-evaluator replication rule is strict: a result is called
replicated only when the direction is compatible and the Holm-adjusted test is
significant for both evaluators. Risk metrics are oriented so lower values favor
\mthree{}.

The directional protocol is intentionally role-aware: evaluators see the reference
and transformation labels because creative-core preservation and semantic drift cannot
be assessed without that relationship. Directional prototype choice therefore measures
traceability-oriented recommendation, not neutral A/B preference.

\section{Results}
\subsection{All individual metrics in the primary study}
Table~\ref{tab:individualall} reports every primary-study individual metric, not only
the significant findings. Technical plausibility is the only result satisfying the
pre-specified cross-evaluator replication criterion. Gemini additionally reports
M3-favorable changes in MVP clarity, unsupported-claim risk, ethical/privacy risk,
overall MVP usefulness, and prototype willingness; these remain evaluator-specific
secondary findings. The full statistics, bootstrap intervals, and tie counts are in
Supplementary Table S1 and the machine-readable file
\texttt{results/tables/S01\_Gemini31\_P1\_Individual\_All\_Metrics.csv}.

\begin{table}[htbp]\centering\scriptsize
\caption{Gemini 3.1 Flash Lite P1: all individual-card metrics ($n=60$ pairs per evaluator). $\Delta=$M3--M1. ``Rep.'' requires Holm significance for both evaluators in the same direction.}
\label{tab:individualall}
\begin{tabular}{lrrrrr}
\toprule
Metric & GPT $\Delta$ & GPT $p_H$ & Gemini $\Delta$ & Gemini $p_H$ & Rep.\\
\midrule
Novelty & -0.067 & 1.000 & -0.217 & .010 & No \\
Conceptual Coherence & +0.042 & 1.000 & +0.067 & .656 & No \\
User Pain Plausibility & +0.000 & — & +0.033 & 1.000 & No \\
Technical Plausibility & +0.150 & .025 & +0.433 & <.001 & Yes \\
MVP Clarity & -0.025 & 1.000 & +0.583 & <.001 & No \\
Business Model Plausibility & +0.133 & .163 & -0.033 & 1.000 & No \\
Unsupported Claim Risk & -0.192 & .110 & -0.333 & <.001 & No \\
Ethical Privacy Risk & -0.033 & 1.000 & -0.117 & .062 & No \\
Overall MVP Usefulness & +0.050 & .089 & +0.233 & .014 & No \\
Prototype Willingness & +0.055 & .110 & +0.367 & <.001 & No \\
\bottomrule
\end{tabular}
\end{table}

\subsection{Directional traceability evidence}
Table~\ref{tab:directionalall} reports all five directional metrics. Both evaluators
rated creative-core preservation near 4/5, practical MVP improvement above 4.27/5,
and overclaim-removal quality above 4/5. Semantic-drift risk was below the midpoint
for both evaluators. This is the main evidence about the intended process mechanism:
preservation of a recognizable creative core alongside practical bounding.

\begin{table}[htbp]\centering\small
\caption{Gemini 3.1 Flash Lite P1: role-aware directional metrics ($n=60$ per evaluator).}
\label{tab:directionalall}
\begin{tabular}{lrrrr}
\toprule
Metric & GPT mean & GPT 95\% CI & Gemini mean & Gemini 95\% CI\\
\midrule
Creative Core Preservation & 3.998 & [3.893, 4.105] & 4.030 & [3.868, 4.182] \\
Novelty Retention & 3.803 & [3.692, 3.917] & 3.408 & [3.230, 3.582] \\
Practical MVP Improvement & 4.370 & [4.325, 4.415] & 4.277 & [4.182, 4.370] \\
Overclaim Removal Quality & 4.100 & [4.032, 4.168] & 4.077 & [3.962, 4.188] \\
Semantic Drift Risk & 2.428 & [2.273, 2.580] & 2.013 & [1.850, 2.187] \\
\bottomrule
\end{tabular}
\end{table}

\begin{table}[htbp]\centering\small
\caption{Gemini 3.1 Flash Lite P1: role-aware directional prototype choice. This is not a neutral preference test.}
\label{tab:prototype}
\begin{tabular}{lrrrrr}
\toprule
Evaluator & Transformation & Reference & Tie & Transform. share excl. ties & Exact $p$\\
\midrule
ChatGPT Plus 5.5 Web & 48 & 12 & 0 & .800 & $3.18\times10^{-6}$\\
Gemini 3.1 Pro Web & 37 & 5 & 18 & .881 & $4.43\times10^{-7}$\\
\bottomrule
\end{tabular}
\end{table}

\subsection{Neutral paired evidence is retained but not used for efficacy inference}
The repository retains detailed neutral paired tables for transparency. Their
evidential status is reported in Table~\ref{tab:status}. Gemini 3.1 P1 neutral
outputs are quarantined from efficacy/equivalence inference because later diagnostics
identified position-sensitive behavior; Gemini neutral v1/v2/v3 outputs are also
archived as excluded preference-consistency artifacts. Gemini 2.5 P1 neutral results
are reported in the supplement but are non-confirmatory after Holm correction.
Consequently, no main paper claim relies on neutral paired results.

\begin{table}[htbp]\centering\small
\caption{Evidence disposition.}
\label{tab:status}
\begin{tabularx}{\textwidth}{@{}lYY@{}}
\toprule
Evidence stream & Repository treatment & Manuscript use\\
\midrule
G3.1 individual & Fully reported, strict schema audits pass & Primary individual inference\\
G3.1 directional & Fully reported, strict schema audits pass & Primary traceability inference\\
G3.1 neutral & Full archive retained; quarantined & No efficacy/equivalence claim\\
G2.5 P0/P1 & Full detailed supplement retained & Supporting/developmental only\\
Human evaluation & Protocol prepared, no outcomes collected & No claim\\
\bottomrule
\end{tabularx}
\end{table}

\section{Discussion}
The evidence supports a limited interpretation of \chvd{}. A speculative concept can
serve as a source of a distinctive creative mechanism. A controlled transformation
can then express that mechanism through a smaller, more testable MVP hypothesis. The
most reliable individual-card effect in the primary corpus is higher
evaluator-assessed technical plausibility. The strongest process evidence is
directional: substantial creative-core preservation, high practical-MVP improvement,
high overclaim-removal quality, and lower semantic-drift risk.

The results indicate a structured novelty--practicalization trade-off rather than an
unconditional quality gain. The study does not establish factual truth, implementation
success, user adoption, profitability, or market viability.

\section{Limitations and reproducibility}
All reported outcomes are LLM-evaluator assessments of written concept cards. They
are not human or expert judgments and are not implementation outcomes. A blinded human
evaluation is planned as external validation. The repository includes this manuscript,
its compiled PDF, all aggregate result tables, source outputs, validators, hashes,
and explicit exclusions. It does not release private A/B keys, candidate-level
unblinding maps, credentials, account/session metadata, or human-participant data.

\section{Conclusion}
Within the evaluated Gemini 3.1 Flash Lite P1 matched corpus, \chvd{} provides
auditable evidence for a constrained transformation from selected speculative
mechanisms to evaluator-assessed testable MVP hypotheses. The findings do not claim
startup readiness. They support a narrower claim of controlled LLM concept
distillation, with human evaluation as the next external validation layer.

\begin{thebibliography}{9}
\bibitem{ji2022survey} Z. Ji et al., ``Survey of Hallucination in Natural Language Generation,'' \textit{ACM Computing Surveys}, 2023.
\bibitem{wang2023spp} Z. Wang et al., ``Unleashing the Emergent Cognitive Synergy in Large Language Models: A Task-Solving Agent through Multi-Persona Self-Collaboration,'' arXiv:2307.05300, 2023.
\bibitem{liu2023geval} Y. Liu et al., ``G-Eval: NLG Evaluation using GPT-4 with Better Human Alignment,'' arXiv:2303.16634, 2023.
\bibitem{zheng2023judge} L. Zheng et al., ``Judging LLM-as-a-Judge with MT-Bench and Chatbot Arena,'' arXiv:2306.05685, 2023.
\end{thebibliography}
\end{document}
